% Template LaTeX initial pour le rapport
% Remplacer progressivement le contenu par la transcription structurée du .docx
\documentclass[11pt,a4paper]{report}
\usepackage[utf8]{inputenc}
\usepackage[T1]{fontenc}
\usepackage[french]{babel}
\usepackage{lmodern}
\usepackage{geometry}
\usepackage{graphicx}
\usepackage{hyperref}
\usepackage{xcolor}
\usepackage{amsmath,amssymb}
\usepackage{booktabs}
\usepackage{enumitem}
\usepackage{longtable}
\usepackage{tabularx}
\usepackage{caption}
\usepackage{float}
\usepackage{listings}
\usepackage{csquotes}
\usepackage{pdfpages}
\geometry{margin=2.5cm}
\hypersetup{colorlinks=true,linkcolor=blue,urlcolor=blue,citecolor=blue}

% Définition d'un environnement pour les objectifs
\newcommand{\objectif}[1]{\paragraph{Objectif}#1}

% Tableaux flexibles
\newcolumntype{Y}{>{\raggedright\arraybackslash}X}

\title{Rapport de Pré-Études Photovoltaïques\\\large Intégration Estimation, Finance et Export}
\author{Entreprise / Auteur}
\date{\today}

\begin{document}
\maketitle
\tableofcontents
\chapter{Introduction}
Ce document formalise le rapport précédemment rédigé dans un fichier Word (Rapport\_EAubonne.docx) en le convertissant vers une structure LaTeX maintenable.

\section{Contexte}
% Décrire le contexte métier / besoin initial

\section{Objectifs}
Les objectifs spécifiques du projet sont :
\begin{enumerate}
  \item Concevoir, intégrer et analyser les données d’estimation de production (Grid, Off-Grid, Tracker).
  \item Réduire le temps de préparation d’un rapport d’estimation.
  \item Diminuer les risques d’erreurs manuelles dans la consolidation multi-outils.
  \item Modéliser les données et implémenter les indicateurs financiers (VAN, TRI, etc.).
  \item Créer une interface ergonomique avec validation des entrées et contrôleurs fiables.
  \item Industrialiser la qualité (tests, couverture, analyse statique, documentation).
\end{enumerate}

\chapter{Architecture}
\section{Vue Composants}
Inclure le diagramme de composants (PlantUML exporté en PNG/PDF) :
\begin{figure}[H]
  \centering
  % \includegraphics[width=0.9\textwidth]{../uml/component-architecture.png}
  \caption{Architecture logique par composants}
\end{figure}

\section{Cas d'utilisation principaux}
\subsection{Estimation PV Grid}
% Préconditions / Scénario nominal / Post-conditions / Alternatifs
\subsection{Graphes Financiers}
% Préconditions / Scénario nominal / Post-conditions / Alternatifs
\subsection{Export PDF / CSV}
% Préconditions / Scénario nominal / Post-conditions / Alternatifs

\chapter{Modélisation}
\section{Données et Domaines}
% Décrire modele.pvgis.*, modele.finance.*
\section{Patrons de conception}
% Strategy (Export), Façade (PVGIS/Export), Builder (PDF)

\chapter{Implémentation}
\section{Technologies}
Java 21, Maven, Swing, HttpClient, JSON (org.json), XChart, PDFBox, Apache POI.
\section{Qualité et Tests}
% Résumé des tests vitrine

\chapter{Résultats et Évaluation}
\section{Tableau Objectifs / Résultats / Méthodes}
\begin{longtable}{p{0.30\textwidth} p{0.36\textwidth} p{0.30\textwidth}}
\toprule
\textbf{Objectif} & \textbf{Résultats (atteints / attendus)} & \textbf{Méthodes d'évaluation} \\
\midrule
Concevoir, intégrer et analyser l’estimation (Grid, Off-Grid, Tracker) & Services unifiés, parseur robuste, normalisation unités & Revue architecture; tests parser; comparaison PVGIS officielle \\
Réduire le temps de préparation d’un rapport & Export PDF/CSV auto; réduction temps manuel (X→Y min) & Mesures chronométrées; feedback utilisateurs; comptage étapes \\
Réduire les erreurs de consolidation & Centralisation calculs + validation; formats homogènes & Historique erreurs vs incidents; tests validation; contrôle exports \\
Implémenter indicateurs financiers & FinancialCalculator (VAN, TRI, payback); conversions validées & Tests formules; comparaison Excel; revue arrondis \\
Interface ergonomique et contrôleurs fiables & UI réactive; états explicites; blocage actions en erreur & Tests contrôleurs; inspection ergonomie; logs erreurs \\
Industrialiser la qualité & Tests ciblés; profil JaCoCo; Sonar A/A; docs UML + README & Rapport JaCoCo; tableau de bord Sonar; revue documentaire \\
\bottomrule
\end{longtable}

\chapter{Conclusion}
% Synthèse finale

\appendix
\chapter{Annexes}
% Ajouter exports de diagrammes, extraits de code, etc.

\end{document}
